% !TEX root = ../main/main.tex
\section{Unified Baseline Comparison}
\label{sec:unified-baseline-comparison}

This section synthesizes performance across all baseline models developed in this chapter, demonstrating that ensemble integration outperforms individual models through complementary strengths.

\subsection{Model Portfolio Summary}

Table~\ref{tab:unified-baseline-summary} presents a comprehensive comparison of all baseline models on the 2024 holdout set (285 games).

\begin{table}[htbp]
\centering
\caption[Unified baseline model comparison]{Unified baseline model comparison on 2024 holdout (285 games). Metrics include spread MAE, win probability AUC, ATS accuracy, calibration (ECE), inference time, and trainable parameters. Bold indicates best-in-class performance.}
\label{tab:unified-baseline-summary}
\begin{adjustbox}{max width=\textwidth}
\small
\begin{tabular}{@{}lcccccc@{}}
\toprule
\textbf{Model} & \textbf{Spread MAE} & \textbf{Win AUC} & \textbf{ATS Acc.} & \textbf{ECE} & \textbf{Inference} & \textbf{Params} \\
\midrule
\textit{Classical Baselines} & & & & & & \\
State-Space (Kalman) & 10.8 & 0.584 & 51.2\% & 0.042 & $<1$ms & 33 \\
Logistic/Probit GLM & 11.2 & 0.572 & 50.8\% & 0.038 & $<1$ms & 12 \\
\midrule
\textit{Machine Learning} & & & & & & \\
XGBoost (293 features) & 10.5 & 0.628 & 52.3\% & 0.051 & 8ms & 2.4K \\
GameOutcomeModel & \textbf{10.3} & \textbf{0.648} & 53.1\% & 0.047 & 50ms & 2.8K \\
\midrule
\textit{Failed Approaches} & & & & & & \\
5-Component Ensemble & 8.3$^*$ & 0.715$^*$ & 62.1\%$^*$ & 0.089$^*$ & 250ms & 12K \\
GNN (Optuna Best) & 12.1 & 0.523 & 45.4\% & 0.124 & 180ms & 340K \\
XGBoost v4 Ridge & 59.1$^\dagger$ & 0.412$^\dagger$ & 38.2\%$^\dagger$ & -- & 85ms & 3.1K \\
\midrule
\textbf{Ensemble (v6.0)} & \textbf{10.1} & 0.641 & \textbf{54.1\%} & \textbf{0.035} & 105ms & 5.2K \\
\midrule
\textit{Theoretical Ceiling} & & & & & & \\
Market (Closing Line) & 9.7 & 0.687 & 50.0\% & 0.018 & -- & -- \\
\bottomrule
\end{tabular}
\end{adjustbox}
\begin{tablenotes}
\footnotesize
\item $^*$Overfit: 70\% validation-to-test Brier degradation. $^\dagger$Target leakage: invalidated (Appendix~\ref{app:orthogonalization}).
\item \textbf{Metrics}: Spread MAE in points. Win AUC = area under ROC curve for win probability. ATS = against-the-spread accuracy. ECE = expected calibration error. Inference = p95 latency on M4 Mac.
\end{tablenotes}
\end{table}

\subsection{Key Findings}

\paragraph{1. Ensemble Integration Outperforms Individual Models.}
The v6.0 ensemble achieves MAE 10.1 (vs.\ 10.3 for best single model) and 54.1\% ATS (vs.\ 53.1\% for GameOutcomeModel), demonstrating that combining complementary models provides measurable improvement. The ensemble reaches 96\% of theoretical market performance (10.1 vs.\ 9.7 MAE).

\paragraph{2. Simple Models Remain Competitive.}
The state-space model (33 parameters, $<1$ms inference) achieves 10.8 MAE---only 7\% worse than the 293-feature GameOutcomeModel that requires 50$\times$ more computation. This validates the principle that \textbf{simpler models generalize better with limited data}.

\paragraph{3. Complexity Without Discipline Fails.}
Three approaches that prioritized model complexity over validation rigor failed dramatically:
\begin{itemize}[itemsep=2pt,parsep=0pt]
    \item \textbf{5-Component Ensemble}: 70\% validation-to-test degradation (overfitting)
    \item \textbf{GNN}: Below-random 45.4\% accuracy (generalization failure)
    \item \textbf{XGBoost v4 Ridge}: 47\% worse than baseline (target leakage)
\end{itemize}

\paragraph{4. Calibration Matters More Than Sharpness.}
The ensemble's superior ECE (0.035 vs.\ 0.047--0.089) translates directly to economic value: well-calibrated probabilities enable Kelly criterion bet sizing that maximizes expected log wealth. Overconfident models (high ECE) lead to bankroll ruin.

\paragraph{5. Market Efficiency Imposes Hard Ceiling.}
No model exceeds market MAE by more than 4\%, confirming that \textbf{vigorish barrier} limits profitability. The closing line incorporates all public information; our models provide marginal edge (if any) only in early-week predictions before market convergence.

\subsection{Variance Decomposition}

The ensemble's improvement comes from diversification across uncorrelated error sources. Table~\ref{tab:variance-decomposition} shows pairwise correlation of model residuals.

\begin{table}[htbp]
\centering
\caption[Model residual correlations]{Pairwise correlation of prediction residuals on 2024 holdout. Lower correlation indicates higher diversification benefit from ensembling.}
\label{tab:variance-decomposition}
\small
\begin{tabular}{@{}lccc@{}}
\toprule
& \textbf{State-Space} & \textbf{XGBoost} & \textbf{GameOutcome} \\
\midrule
State-Space & 1.00 & 0.72 & 0.68 \\
XGBoost & 0.72 & 1.00 & 0.84 \\
GameOutcomeModel & 0.68 & 0.84 & 1.00 \\
\bottomrule
\end{tabular}
\end{table}

The moderate correlations (0.68--0.84) confirm that models capture different aspects of game outcomes, justifying ensemble combination. State-space and GameOutcomeModel have lowest correlation (0.68), explaining why their combination yields the largest improvement.

\subsection{Computational Trade-offs}

\paragraph{Accuracy-Latency Pareto Frontier.}\label{fig:pareto-baseline}
The state-space model (10.8 MAE, $<1$ms) anchors the efficient frontier with best accuracy-per-millisecond. For latency-constrained applications (live betting), state-space dominates. For batch predictions where 105ms latency is acceptable, the v6.0 ensemble provides optimal accuracy. GameOutcomeModel and ensemble dominate for accuracy-prioritized deployments.

\subsection{Implications for Production}

Based on this analysis, we deploy a tiered prediction strategy:

\begin{enumerate}[itemsep=2pt,parsep=0pt]
    \item \textbf{Pre-game (batch)}: v6.0 ensemble with full 293 features, 105ms latency acceptable
    \item \textbf{Live (real-time)}: State-space model only, $<1$ms latency required
    \item \textbf{Early-week exploration}: GameOutcomeModel for market inefficiency detection before line stabilization
\end{enumerate}

The unified comparison demonstrates that \textbf{model diversity + rigorous validation} trumps sophisticated individual models. The ensemble's 54.1\% ATS represents the production ceiling after accounting for vigorish and market efficiency constraints.
